\begin{abstract}
Solar power and food production increasingly compete for the same land, while warming exposes crops to more heat and drought. Agrivoltaic systems, which grow crops under photovoltaic (PV) modules, can relieve this conflict, but how much electricity farmland can produce without compromising food security depends on how sunlight is shared between modules and plants, day by day, for every crop and climate. We propose SALS (stress-aware light sharing), an algorithm that learns the design and the daily control of single-axis tracker agrivoltaics end-to-end through a differentiable simulator. The simulator couples solar geometry, backtracking trackers, PV output, crop-level light and microclimate, a soil-water balance and the SIMPLE crop model with heat and drought stress. A design network maps site descriptors to the ground coverage ratio, and a control network maps the daily crop, soil and weather state to the tracker's light-sharing level. Both are trained by back-propagation through whole seasons to maximise electricity subject to a 90\% yield-retention floor. We applied SALS to 400 sites representing the global harvested area of wheat, rice, maize, soybean and potato, with NASA POWER weather for 2001--2019 and warming of +1.5, +2 and +3\,\textdegree C. On geographically held-out sites and years, SALS produced 1019\,MWh\,ha$^{-1}$\,yr$^{-1}$ (72\% of a conventional PV plant) at a mean yield retention of 0.91. This gave a land-equivalent ratio of 1.63, compared with 1.20 for static agrivoltaic designs and 1.42 for an expert stress rule ($p<10^{-4}$). Its advantage grew with warming, to an LER of 1.75 at +3\,\textdegree C, as the learned policy shifted from sharing light during canopy development to shading crops during heat. The yield floor was met on average but in only 45--49\% of individual seasons, which calls for chance-constrained extensions. Scaled to 5\% of the harvested area of the five crops, learned agrivoltaics has a technical potential of about 31{,}000\,TWh\,yr$^{-1}$, comparable to current global electricity generation.
\end{abstract}

\begin{keyword}
Agrivoltaics \sep Differentiable simulation \sep Food--energy nexus \sep Climate adaptation \sep Solar tracking \sep Crop modelling \sep Physics-informed machine learning
\end{keyword}
